\subsection{Recognizing relations over the prime field}
Geometric integrality permits degree reductions beyond those implied
by radicality. An algebraic element of the function field must already
belong to the prime field. A linearized polynomial relation can
therefore force a hidden factor to vanish, provided its coefficient
sum is nonzero. The following theorem recognizes such a relation with
an unknown number of lower Frobenius blocks and unknown coefficients.
Additive-polynomial composition and its skew-polynomial interpretation
are classical. In particular, \citet[\S6, Theorem~6.1, pp.~22--24 of
the author version]{sp-giesbrecht} gives deterministic decomposition
in time polynomial in the Frobenius degree at fixed characteristic,
including characteristic two. That input already specifies a
univariate additive polynomial. Here the sparse homogeneous equation
hides both its additive coefficients and its multivariate argument
$U/V$; the exponent gap below permits their recovery.

\begin{theorem}[Sparse recognition of a linearized relation]
\label{sp:thm:linearized-preprocessing}
Fix a prime $p$. Let a homogeneous sparse list with binary exponents
generate $K\subset S=\Fp[x_0,\ldots,x_N]$, and suppose that
$I=K^{\mathrm{sat}}$ has geometrically integral quotient $S/I$.
Consider a nonzero row admitting the exact identity
\begin{equation}
 F=U^Q+\sum_{i=0}^{s}a_iV^{Q-p^i}U^{p^i},
 \qquad Q=p^e,\quad 0\le s<e,
 \label{sp:eq:linearized-relation}
\end{equation}
where $U\ne0$ is homogeneous of degree $d\ge1$, $V$ is any
coefficient-one monomial of degree $d$, and $a_i\in\Fp$. Assume
\begin{equation}
 (d+1)p^s\le Q,\qquad
 \lambda:=1+\sum_{i=0}^{s}a_i\ne0.
 \label{sp:eq:linearized-gap}
\end{equation}
None of $U,V,e,s,d$ or the coefficient vector is given. The number
$s$ is not fixed, and zero lower coefficients are allowed.

There is a deterministic polynomial sparse-bit algorithm that finds
and verifies such an identity whenever one exists and replaces $F$
by $U$. Each accepted replacement satisfies
\[
 K\subseteq K'\subseteq I,\qquad (K')^{\mathrm{sat}}=I,
\]
introduces no coordinate index, does not increase the number of
terms, and divides the numerical degree by $Q\ge p$.
The algorithm may choose a minimum-degree replacement among its
passing candidates. Repeated use, interleaved with rank-one digit
reduction and batched valid coordinate-content division, terminates
in polynomial sparse-bit time. An ordinary whole Frobenius power
may also be handled by its existing root-extraction rule.
\end{theorem}
\begin{proof}
Put $R=S/I$ and $E=\operatorname{Frac}(R)$. If $z\in E$ is
algebraic over $\Fp$, then $\Fp(z)=\mathbb F_{p^b}$ for some $b$.
Tensoring its inclusion in $E$ with $\overline{\Fp}$ preserves
injectivity. For $b>1$ the source is a product of fields and has
nonzero zero divisors, whereas the target is a localization of the
domain $R\otimes_{\Fp}\overline{\Fp}$. Hence $b=1$.
This is the usual geometric-integrality condition
\citep[Definition~10.49.1 and Lemma~10.49.3]{sp-stacks-integral}.
The saturated homogeneous coordinate ring of a geometrically
integral projective scheme has this property: after base change,
its punctured affine cone is a $\mathbb G_m$-torsor over the
integral scheme, and its affine closure is integral.

Reduce \eqref{sp:eq:linearized-relation} modulo $I$.
If $V=0$, then $U^Q=0$ and $U=0$ in the domain $R$.
Otherwise $z=U/V\in E$ satisfies
\[
 z^Q+\sum_{i=0}^{s}a_i z^{p^i}=0.
\]
It is algebraic over $\Fp$ and therefore lies in $\Fp$.
The equation becomes $\lambda z=0$, forcing $z=0$.
Thus $U\in I$. Every term of $F$ is divisible by $U$,
which proves the ideal sandwich and equality of saturations.
This safety argument does not require the gap or monomial hypotheses;
those hypotheses will make the unknown relation recognizable.

Write $D=\deg F$ and $T=|\operatorname{supp}F|$.
Enumerate $Q=p^e\le D$ with $Q\mid D$, set $d=D/Q$, and
enumerate $0\le s<e$ satisfying \eqref{sp:eq:linearized-gap}'s
degree inequality. There are $O(\log^2(D+1))$ pairs $(e,s)$.
For each pair, take the \emph{$Q$-section} $W$ of $F$: select
the terms whose exponent vectors are coordinatewise divisible by
$Q$, and divide those vectors by $Q$, retaining their coefficients.
Reject a zero or wrong-degree section.

To determine what a true relation gives, write
$U=\sum_u c_u x^u$ and $V=x^v$, with $|u|=|v|=d$.
A term of its $i$th lower block has exponent
\[
 (Q-p^i)v+p^iu=Qv+p^i(u-v).
\]
Every coordinate of $p^i(u-v)$ has absolute value at most
$dp^s\le Q-p^s<Q$. It is coordinatewise divisible by $Q$
only when $u=v$. If $c_v=0$ when $V$ is absent from $U$, then
\begin{equation}
 W=U+(\lambda-1)c_vV.
 \label{sp:eq:leading-section}
\end{equation}
Because $\lambda\ne0$, $W$ has exactly the support of $U$.
No membership in $I$ is inferred from taking a section alone.

First suppose $V$ occurs in this support. Enumerate its at most
$T$ possibilities and the $p-1$ possible nonzero values of
$\lambda$. Leave all other coefficients of $W$ unchanged and
divide its $V$ coefficient by $\lambda$ to propose $U$.
For these fixed candidates form the sparse polynomials
\[
 B_i=V^{Q-p^i}U^{p^i}\quad(0\le i\le s),\qquad H=F-U^Q.
\]
Solve over $\Fp$ the coefficient equations
\begin{equation}
 H=\sum_{i=0}^{s}a_i B_i,
 \qquad \sum_{i=0}^{s}a_i=\lambda-1.
 \label{sp:eq:linearized-system}
\end{equation}
Use the union of the supports of $H$ and the $B_i$ as the row
indices of this linear system. Its unknowns are only the $s+1$
scalars $a_i$. A consistent system gives one coefficient vector;
there is no enumeration of $p^{s+1}$ vectors. Verify the complete
identity and every degree condition before accepting. The true
support-present representation is among the candidates.

For the other case set $U=W$ and again put $H=F-U^Q$.
If $H=0$, ordinary whole-power extraction applies. It also admits
a relation with all lower coefficients zero, $s=0$, and any
support monomial of $U$ as $V$, when $d+1\le Q$.
Otherwise select a surviving exponent vector $\alpha$ of $H$
and put $q=p^s$, $M=Q-q$. In a true representation it has the form
\[
 \alpha=Mv+\beta,\qquad
 \beta=(q-p^i)v+p^iu\ge0,\qquad |\beta|=qd\le M.
\]
Cancellation may delete terms but cannot produce a new exponent,
so this description holds for any chosen surviving term.
Consequently $\lfloor\alpha/M\rfloor$ equals either $v$ or
$v+e_j$ for one coordinate unit vector $e_j$. The latter requires
$qd=M$ and $\beta=Me_j$. Enumerate that quotient and the vectors
obtained by subtracting one coordinate unit, retaining only
nonnegative vectors of total degree $d$. There are at most $v_0+1$
candidates, where $v_0$ is the number of active coordinates in
the input row. For each candidate $V$ and each
$\lambda\in\Fp^*$, solve \eqref{sp:eq:linearized-system} and
verify the full identity. This includes the boundary case of
equality in the degree gap and the true absent-support representation.

The searches have polynomial sparse-bit cost. Each $B_i$ is obtained
by multiplying exponent integers and adding a monomial shift;
its support has at most $T$ terms. The coefficient matrix has at
most $(s+2)T+T+1$ rows and $s+1\le\log_p D$ columns.
Gaussian elimination, sparse collection, and exact verification
are polynomial in the written input length. Intermediate exponent
integers have polynomial bit length, since a candidate identity is
homogeneous of degree $D$. Each positive coordinate exponent of
$V$ in the nonzero-remainder case occurs in that remainder, because
$M>0$. Thus absent coordinate identifiers need never be enumerated.

Every accepted $U$ is a section with at most one coefficient changed,
so it has no more than $T$ terms and introduces no coordinate.
All passing candidate outputs can be compared by degree, with a
fixed ordering to break ties. Each new step decreases degree by
at least a factor two. The known digit operation decreases it by
at least the fraction $1/(v_0+2)$ and also does not increase
support or active coordinates. Batched eligible monomial content
is removed at most once between such steps. Hence there are
$O(v_0\log(D_0+1))$ successful degree-reducing steps from initial
degree $D_0$, proving polynomial termination and the combined cost.
\end{proof}
